\documentclass[11pt]{article}
\usepackage{times}
\usepackage{fullpage}
\usepackage{url}
\usepackage{hyperref}
\usepackage{fancyhdr}
\usepackage{graphicx}
\usepackage{enumitem}
\usepackage{subcaption}
\usepackage{amsmath, amsfonts, amsthm, fouriernc}
\usepackage{IEEEtrantools}
\usepackage{parskip}

\pagestyle{fancy}
\addtolength{\headheight}{2em}
\addtolength{\headsep}{1.5em}
\lhead{\large{ME 2801}}

\usepackage{sectsty}
\sectionfont{\fontsize{12}{8}\selectfont}

\begin{document}

\begin{center}
  \Large{\bf{Lab Assignment: USV System Identification}}
\end{center}

This assignment accompanies the USV open-loop step-response experiments.
You will process flight-log data, identify first-order plant models for surge and
yaw, and validate the models against measured data.

\textbf{Deliverable.}  A brief lab report (PDF) with figures and MATLAB code
appended.

\bigskip\hrule\bigskip

% ---------------------------------------------------------------
\section*{1\quad Data Processing}

\begin{enumerate}

\item Convert the ArduPilot \texttt{.BIN} log to \texttt{.mat} format using
  \texttt{bin2mat.py}.  Run \texttt{proc\_openloop\_response.m} to extract and
  plot the surge and yaw step responses.

\item Identify the start and end of each step (the interval where the actuator
  command is approximately constant).  Trim the time series to the relevant
  windows and shift time so that $t = 0$ at the start of each step.

\item Plot the trimmed surge (speed vs.\ time) and yaw (yaw rate vs.\ time)
  responses on separate axes.  Label axes with units.

\end{enumerate}

% ---------------------------------------------------------------
\section*{2\quad Surge Model Identification}

The surge plant is modeled as a first-order system:
\[
  G_{\mathrm{surge}}(s) = \frac{K_{\mathrm{surge}}}{\tau_{\mathrm{surge}}\,s + 1}
\]
where the input is normalized throttle ($\pm 1$) and the output is forward speed (m/s).

\begin{enumerate}[resume]

\item From the step-response data, estimate the steady-state speed $y_{ss}$ and
  the normalized throttle command $u_0$.  Compute the DC gain:
  \[
    K_{\mathrm{surge}} = \frac{y_{ss}}{u_0}
  \]

\item Estimate the time constant $\tau_{\mathrm{surge}}$ from the response.
  (The output reaches $63.2\%$ of $y_{ss}$ at $t = \tau$.)

\item Construct the identified transfer function as a MATLAB \texttt{tf} object.
  Simulate the model's step response and overlay it on the measured data.
  Comment on the fit quality.

\end{enumerate}

% ---------------------------------------------------------------
\section*{3\quad Yaw Model Identification}

The yaw plant is modeled as an integrator with damping:
\[
  G_{\mathrm{yaw}}(s) = \frac{K_{\mathrm{yaw}}}{s\,(\tau_{\mathrm{yaw}}\,s + 1)}
\]
where the input is normalized rudder command ($\pm 1$) and the output is yaw rate (rad/s).

\begin{enumerate}[resume]

\item From the step-response data, estimate the steady-state yaw rate $r_{ss}$
  and the normalized rudder command $\delta_0$.  Compute the gain parameter
  $K_{\mathrm{yaw}} = r_{ss}/\delta_0$.

\item Estimate $\tau_{\mathrm{yaw}}$ from the initial transient shape.

\item Construct the identified transfer function.  Simulate and overlay on the
  measured yaw rate.  Comment on the fit.

\end{enumerate}

% ---------------------------------------------------------------
\section*{4\quad Discussion}

\begin{enumerate}[resume]

\item What are the dominant sources of uncertainty in this identification
  procedure?  How would you reduce them in a follow-on experiment?

\item How would wind, waves, or current affect the identified parameters?

\end{enumerate}

\end{document}
